% Deux matrices d'attention sur les mêmes jetons.
%   gauche : attention plate — un coefficient libre par paire (i,j)
%   droite : attention hiérarchique — la hiérarchie {A,B,C} découpe la matrice
%            en blocs, et toutes les paires entre deux sous-arbres frères
%            partagent une seule valeur (block constraint, HSA, NeurIPS 2025).
% Les neuf tons de droite sont les moyennes des blocs de gauche : la contrainte
% de blocs est bien une *compression* de la matrice plate, jamais un ajout.
\begin{tikzpicture}[xscale=1.0, yscale=1.0]
    \useasboundingbox (-0.75, -1.05) rectangle (8.55, 3.35);

    \tikzset{cadre/.style={draw=gray!50, line width=0.9pt}}
    \tikzset{titre/.style={font=\small\bfseries, anchor=south, align=center}}
    \tikzset{sous/.style={font=\scriptsize, anchor=north, align=center, text=gris_fonce_inria}}
    \tikzset{lettre/.style={font=\scriptsize\bfseries, text=inria-rouge}}

    \def\cs{0.32}   % côté d'une cellule

    % ===================== gauche : attention plate =====================
    \begin{scope}[shift={(0,0)}]
        \node[titre, text=gris_fonce_inria] at (1.28,2.98) {attention plate};
        \foreach \i in {0,...,7} {
            \foreach \j in {0,...,7} {
                \pgfmathsetmacro{\gg}{max(6, min(92,
                    12 + 78*exp(-(\i-\j)*(\i-\j)/4) + 9*sin((2.7*\i+1.3*\j)*57.2957795)))}
                \fill[inria-2024-bleu-canard!\gg] (\i*\cs,\j*\cs) rectangle ++(\cs,\cs);
            }
        }
        \draw[cadre] (0,0) rectangle (8*\cs,8*\cs);
        \node[sous] at (1.28,-0.16) {$N^2$ coefficients \textbf{libres} :\\chaque paire $(i,j)$ est jugée seule};
    \end{scope}

    % ===================== droite : attention hiérarchique ==============
    \begin{scope}[shift={(5.44,0)}]
        \node[titre, text=inria-rouge] at (1.28,2.98) {attention hiérarchique};

        % blocs : une seule valeur par couple de sous-arbres frères
        \def\bornes{{0,3,5,8}}
        \foreach \a/\b/\gg in {%
            0/0/72, 0/1/33, 0/2/15,
            1/0/34, 1/1/83, 1/2/34,
            2/0/15, 2/1/33, 2/2/72} {
            \pgfmathsetmacro{\xa}{\bornes[\a]*\cs}
            \pgfmathsetmacro{\xb}{\bornes[\a+1]*\cs}
            \pgfmathsetmacro{\ya}{\bornes[\b]*\cs}
            \pgfmathsetmacro{\yb}{\bornes[\b+1]*\cs}
            \fill[inria-2024-bleu-canard!\gg] (\xa,\ya) rectangle (\xb,\yb);
        }
        \draw[cadre] (0,0) rectangle (8*\cs,8*\cs);
        \foreach \k in {3,5} {
            \draw[inria-rouge, line width=1.0pt] (\k*\cs,0) -- (\k*\cs,8*\cs);
            \draw[inria-rouge, line width=1.0pt] (0,\k*\cs) -- (8*\cs,\k*\cs);
        }
        % les trois sous-arbres, lus sur les deux bords
        \foreach \x/\L in {0.48/A, 1.28/B, 2.08/C} {
            \node[lettre, anchor=south] at (\x,2.62) {$\L$};
            \node[lettre, anchor=east]  at (-0.06,\x) {$\L$};
        }
        \node[font=\scriptsize, text=gris_fonce_inria] at (0.48,2.08) {$\theta_{A,C}$};
        \node[font=\scriptsize, text=gris_fonce_inria] at (2.08,0.48) {$\theta_{A,C}$};
        \node[sous] at (1.28,-0.16) {\textbf{un seul} coefficient par couple\\de sous-arbres frères};
    \end{scope}

    % ===================== la flèche du milieu ==========================
    \draw[-{Latex[length=3mm]}, gris_fonce_inria, line width=1.1pt] (2.90,1.28) -- (4.90,1.28);
    \node[font=\scriptsize, text=gris_fonce_inria, anchor=south, align=center]
        at (3.90,1.42) {la hiérarchie\\découpe en blocs};
    \node[font=\scriptsize, text=inria-rouge, anchor=north, align=center]
        at (3.90,1.12) {non plus les \textbf{paires},\\mais les \textbf{nœuds}};
\end{tikzpicture}
